\section*{Chapitre \ref{ch:b2:genfun} --- Fonctions génératrices des probabilités}

\begin{solution}{exo:b2:genfun:1}
Dé équilibré~: $G(t) = \frac16(t + t^2 + \dots + t^6) = \frac t6(1 + t
+ \dots + t^5)$. Si la somme de deux dés équilibrés était uniforme sur
$\{2, \dots, 12\}$, alors
\[
G(t)^2 = \frac{t^2}{36}\,h(t)^2
= \frac{t^2}{11}\sum_{k=0}^{10}t^k ,
\qquad h(t) = 1 + t + \dots + t^5 .
\]
Or $h$ est un polynôme réel de degré impair $5$, donc il a une racine
réelle (théorème des valeurs intermédiaires~; concrètement $h(-1) =
0$), d'où $h^2$ a une racine réelle. Mais $\sum_{k=0}^{10}t^k$ n'en a
aucune~: il est positif pour $t \geq 0$, et pour $t < 0$ il vaut
$\frac{t^{11} - 1}{t - 1}$, quotient de deux nombres négatifs.
Contradiction --- la somme de deux dés équilibrés n'est jamais uniforme
(comme le confirme la \omterm{def:b2:randomvar:law}{loi} triangulaire familière des sommes de dés).
\end{solution}

\begin{solution}{exo:b2:genfun:2}
\emph{Binomiale~:} $G(t) = (1 - p + pt)^n$, $G'(t) = np(1 - p +
pt)^{n-1}$, $G''(t) = n(n-1)p^2(1 - p + pt)^{n-2}$, donc
\[
\E(X) = G'(1) = np,
\qquad
V(X) = G''(1) + G'(1) - G'(1)^2
= n(n-1)p^2 + np - n^2p^2 = np(1-p).
\]
\emph{Géométrique} ($q = 1 - p$)~: $G(t) = \frac{pt}{1 - qt}$, donc
$G'(t) = \frac{p}{(1 - qt)^2}$ et $G''(t) = \frac{2pq}{(1 -
qt)^3}$~; en $t = 1$ (en utilisant $1 - q = p$)~:
\[
\E(X) = \frac{p}{p^2} = \frac1p,
\qquad
V(X) = \frac{2q}{p^2} + \frac1p - \frac{1}{p^2}
= \frac{2q + p - 1}{p^2}
= \frac{q}{p^2} ,
\]
en accord avec l'~\cref{exo:b2:randomvar:1} avec moins de travail.
\end{solution}

\begin{solution}{exo:b2:genfun:3}
Non, même avec des pipages différents. Supposons que $X, Y$ soient des
\omterm{def:b2:randomvar:law}{lois} sur $\{1, \dots, 6\}$ de somme uniforme. Alors $G_X(t) = t\,a(t)$ et
$G_Y(t) = t\,b(t)$ avec $a, b$ des polynômes réels de degré \emph{au
plus} $5$ --- et leurs degrés doivent sommer à $10$ (la somme atteint
$12$ avec probabilité positive), donc $\deg a = \deg b = 5$, tous deux
impairs. Comme dans l'~\cref{exo:b2:genfun:1},
\[
a(t)\,b(t) = \frac{1}{11}\sum_{k=0}^{10}t^k
\]
forcerait une racine réelle à gauche (tout polynôme réel de degré
impair en a une) et aucune à droite. Donc aucun pipage de deux dés
\omterm{def:b2:proba:independence}{indépendants} --- identique ou non --- ne produit une somme uniforme.
\end{solution}

\begin{solution}{exo:b2:genfun:4}
D'après le~\cref{thm:b2:genfun:compound} avec $G_N(s) = e^{\lambda(s-1)}$
et $G_X(t) = 1 - p + pt$~:
\[
G_S(t) = e^{\lambda(1 - p + pt - 1)} = e^{\lambda p(t - 1)} :
\]
$S \sim \mathcal{P}(\lambda p)$. Le nombre rejeté $D = N - S$
compte les mêmes objets conservés avec probabilité $1 - p$, donc par le
même calcul $D \sim \mathcal{P}(\lambda(1 - p))$.
Indépendance, directement~: pour $j, k \in \N$,
\begin{align*}
\P(S = j,\ D = k)
&= \P(N = j + k)\,\binom{j+k}{j}p^jq^k
= e^{-\lambda}\frac{\lambda^{j+k}}{(j+k)!}\,
\frac{(j+k)!}{j!\,k!}\,p^jq^k\\
&= \Bigl(e^{-\lambda p}\frac{(\lambda p)^j}{j!}\Bigr)
\Bigl(e^{-\lambda q}\frac{(\lambda q)^k}{k!}\Bigr)
\end{align*}
avec $q = 1 - p$~: la \omterm{def:b2:randomvar:law}{loi} jointe se factorise en
$\mathcal{P}(\lambda p) \otimes \mathcal{P}(\lambda q)$. Un flux de
Poisson scindé au hasard donne des flux de Poisson \emph{\omterm{def:b2:proba:independence}{indépendants}}
--- un petit miracle constamment utilisé en théorie des files
d'attente.
\end{solution}

\begin{solution}{exo:b2:genfun:5}
Les temps d'attente entre faces consécutives sont des variables
géométriques $\mathcal{G}(p)$ indépendantes (absence de mémoire~:
après chaque face le jeu recommence), donc $T_r = W_1 + \dots + W_r$ et
la multiplicativité (\cref{thm:b2:genfun:product}) donne
\[
G_{T_r}(t) = \Bigl(\frac{pt}{1 - qt}\Bigr)^{r},
\qquad
\E(T_r) = r\,\E(W_1) = \frac rp,
\qquad
V(T_r) = r\,V(W_1) = \frac{rq}{p^2}
\]
($q = 1 - p$~; les \omterm{def:b2:randomvar:variance}{variances} s'ajoutent par indépendance).
Développement~: par la série du binôme généralisée
(\cref{ch:b2:powerseries}), $(1 - qt)^{-r} = \sum_{m\geq0}
\binom{m + r - 1}{r - 1}q^mt^m$, donc le coefficient de $t^n$ dans
$p^rt^r(1 - qt)^{-r}$ est (avec $m = n - r$)
\[
\P(T_r = n) = \binom{n-1}{r-1}p^r(1-p)^{n-r},
\qquad n \geq r ,
\]
la \omterm{def:b2:randomvar:law}{loi} \emph{binomiale négative} --- combinatoirement~: la $r$-ième
face tombe au lancer $n$ si et seulement si les $r - 1$ faces
précédentes choisissent leurs places parmi les $n - 1$ premiers
lancers.
\end{solution}

\begin{solution}{exo:b2:genfun:6}
$m = 1\cdot\frac38 + 2\cdot\frac38 + 3\cdot\frac18 = \frac{3 + 6 +
3}{8} = \frac32 > 1$~: surcritique. La \omterm{ex:b2:powerseries:fibonacci}{fonction génératrice} est
\[
G(t) = \frac{1 + 3t + 3t^2 + t^3}{8} = \frac{(1 + t)^3}{8} ,
\]
donc les points fixes vérifient $(1 + t)^3 = 8t$, c'est-à-dire $t^3 +
3t^2 - 5t + 1 = 0$. En factorisant la racine garantie $t = 1$~:
\[
t^3 + 3t^2 - 5t + 1 = (t - 1)\bigl(t^2 + 4t - 1\bigr),
\]
et $t^2 + 4t - 1 = 0$ donne $t = -2 \pm \sqrt5$. La racine dans
$\intco{0}{1}$ est $\sqrt5 - 2 \approx 0.236$~: d'après le~\cref{thm:b2:genfun:extinction},
\[
q = \sqrt 5 - 2 .
\]
(Une vérification plaisante~: la \omterm{def:b2:randomvar:law}{loi} de reproduction est celle de $3$
pièces équilibrées indépendantes, $Z_1 \sim \mathcal{B}(3, \frac12)$.)
\end{solution}

\begin{solution}{exo:b2:genfun:7}
\emph{\omterm{def:b2:randomvar:expectation}{Espérance}.} D'abord $\E(Z_n) = m^n$~: d'après le~\cref{thm:b2:genfun:compound}, $\E(Z_{n+1}) = \E(Z_n)\,m$, et
$\E(Z_0) = 1$. La famille $\bigl(Z_n(\omega)\P(\{\omega\})
\bigr)_{n, \omega}$ est positive, donc Fubini pour les familles
s'applique sans condition~:
\[
\E(Y) = \sum_{n=0}^{\infty}\E(Z_n)
= \sum_{n=0}^\infty m^n = \frac{1}{1 - m} < \infty
\]
(en particulier $Y$ est presque sûrement finie~: cohérent avec
l'extinction certaine dans le cas sous-critique).

\emph{Équation fonctionnelle.} Décomposons la population selon les
enfants de l'ancêtre~: si l'ancêtre a $Z_1 = k$ enfants, la
descendance totale est $Y = 1 + Y_1 + \dots + Y_k$, où $Y_i$ est la
descendance totale de la lignée du $i$-ième enfant --- et les $Y_i$
sont des copies indépendantes de $Y$, indépendantes de $Z_1$ (des
lignées distinctes utilisent des \omterm{def:b2:proba:space}{événements} de reproduction disjoints
et \omterm{def:b2:proba:independence}{indépendants}). En conditionnant par $Z_1$ comme dans le~\cref{thm:b2:genfun:compound}~:
\[
H(t) = \E\bigl(t^Y\bigr)
= t\sum_{k=0}^\infty \P(Z_1 = k)\,H(t)^k
= t\,G\bigl(H(t)\bigr),
\]
le facteur $t$ rendant compte de l'ancêtre lui-même. (Pour la
\omterm{def:b2:randomvar:law}{loi} $p_0 = 1 - p$, $p_2 = p$ du branchement binaire, cette équation
quadratique en $H$ se résout explicitement et se développe --- les
\omterm{ex:b2:powerseries:catalan}{nombres de Catalan} du~\cref{ch:b2:powerseries} comptent les arbres
généalogiques.)
\end{solution}

\begin{solution}{exo:b2:genfun:8}
Soit $R > 1$ le rayon et fixons $t \in \intoo{1}{R}$. Alors
$\E(t^X) = G(t) < \infty$, et l'inégalité de Markov
(\cref{thm:b2:randomvar:markov}) appliquée à la variable positive
$t^X$ au niveau $t^n$~:
\[
\P(X \geq n) = \P\bigl(t^X \geq t^n\bigr)
\leq \frac{G(t)}{t^n} = C\rho^n,
\qquad C = G(t),\quad \rho = \frac1t \in \intoo{0}{1}.
\]
\emph{Réciproque~:} si $\P(X \geq n) \leq C\rho^n$, alors $p_n \leq
\P(X \geq n) \leq C\rho^n$, donc pour $\abs t < \frac1\rho$ la
série $\sum p_n\abs t^n$ est dominée par la série géométrique
convergente $C\sum(\rho\abs t)^n$~: le rayon est au moins
$\frac1\rho > 1$. Le rayon de la \omterm{ex:b2:powerseries:fibonacci}{fonction génératrice} et la
décroissance géométrique de la queue sont deux faces de la même
propriété.
\end{solution}

\begin{solution}{exo:b2:genfun:9}
Écrivons $p_k^{(n)} = \P(X_n = k)$, $p_k = \P(X = k)$.

\emph{Cas $k = 0$.} Pour $t \in \intoo{0}{1}$ et toute \omterm{def:b2:randomvar:law}{loi}
$(q_j)$ avec $\sum_j q_j \leq 1$~:
\[
\Bigl|\,q_0 - \sum_j q_jt^j\Bigr|
= \sum_{j \geq 1} q_j t^j
\leq \sum_{j\geq1}t^j = \frac{t}{1 - t} .
\]
D'où
\[
\abs{p_0^{(n)} - p_0}
\leq \frac{2t}{1 - t}
+ \abs{G_{X_n}(t) - G_X(t)} .
\]
Étant donné $\varepsilon > 0$, choisir $t$ tel que $\frac{2t}{1-t} <
\frac\varepsilon2$, puis $n_0$ tel que le dernier terme soit $<
\frac\varepsilon2$ pour $n \geq n_0$~: donc $p_0^{(n)} \to p_0$.

\emph{Étape de récurrence.} Supposons $p_j^{(n)} \to p_j$ pour $j < k$.
Considérons les fonctions \emph{décalées}
\[
g_n(t) = \frac{G_{X_n}(t) - p^{(n)}_0}{t}
= \sum_{j\geq0} p^{(n)}_{j+1}t^j,
\qquad
g(t) = \frac{G_X(t) - p_0}{t} ,
\]
\omterm{ex:b2:powerseries:fibonacci}{fonctions génératrices} des suites sous-probabilistes
$(p^{(n)}_{j+1})_j$ (masse totale $\leq 1$, ce qui est tout ce que
l'argument $k = 0$ utilisait). Pour $t \in \intoo{0}{1}$ fixé, $g_n(t)
\to g(t)$ par hypothèse et le cas $k = 0$. Appliquer l'argument $k =
0$ à $g_n$ donne $p_1^{(n)} \to p_1$~; itérer le décalage $k$ fois
donne $p_k^{(n)} \to p_k$ pour tout $k$. (C'est l'instance discrète et
élémentaire du théorème de continuité de Lévy, dont la forme générale
--- pour les fonctions caractéristiques --- est un jalon de la
troisième année.)
\end{solution}

\begin{solution}{exo:b2:genfun:10}
Ponctuellement, $\frac{1 + (-1)^X}{2}$ vaut $1$ quand $X$ est pair
et $0$ quand impair, donc en prenant les \omterm{def:b2:randomvar:expectation}{espérances} (transfert),
\[
\P(X \text{ pair}) = \frac{1 + \E\bigl((-1)^X\bigr)}2 =
\frac{1 + G_X(-1)}2 .
\]
Poisson~: $\frac{1 + \eu^{-2\lambda}}2 \to \frac12$ quand
$\lambda$ croît. Binomiale~: $\frac{1 + (1 - 2p)^n}2$. Dans les deux
cas $G_X(-1) \to 0$ dit que la parité de $X$ devient une pièce
équilibrée~: la \omterm{def:b2:randomvar:law}{loi} s'étale sur de nombreux entiers et oublie sa
parité.
\end{solution}

\begin{solution}{exo:b2:genfun:11}
$t + \dots + t^6 = t\,\frac{1 - t^6}{1 - t}$ et $1 - t^6 =
(1 - t)(1 + t)(1 + t + t^2)(1 - t + t^2)$, donnant la
factorisation annoncée. Pour le premier dé, $(1 + t)(1 + t + t^2) = 1 +
2t + 2t^2 + t^3$, donc $\frac{t(1+t)(1+t+t^2)}6 = \frac{t +
2t^2 + 2t^3 + t^4}6$~: faces $\{1, 2, 2, 3, 3, 4\}$. Pour le
second, en développant
\[
(1 + 2t + 2t^2 + t^3)(1 - t + t^2)^2 = 1 + t^2 + t^3 + t^4 +
t^5 + t^7,
\]
donc $\frac{t(1+t)(1+t+t^2)(1-t+t^2)^2}6 = \frac{t + t^3 + t^4
+ t^5 + t^6 + t^8}6$~: faces $\{1, 3, 4, 5, 6, 8\}$. Le
produit des deux \omterm{ex:b2:powerseries:fibonacci}{fonctions génératrices} regroupe les six
facteurs en $\bigl(\frac{t(1+t)(1+t+t^2)(1-t+t^2)}6
\bigr)^2$, le carré de la fonction du dé standard~: la
paire de Sicherman a exactement la \omterm{def:b2:randomvar:law}{loi} standard pour le total ---
les \omterm{ex:b2:powerseries:fibonacci}{fonctions génératrices} classent tous ces regroupements.
\end{solution}

\begin{solution}{exo:b2:genfun:12}
Soient $A$ et $B$ les \omterm{ex:b2:powerseries:fibonacci}{fonctions génératrices} de la durée restante
partant de «~aucune face en cours~» et «~une face en cours~». Un
lancer est dépensé, puis~: depuis l'état $0$, pile revient à l'état
$0$, face passe à l'état $1$~; depuis l'état
$1$, face termine le jeu, pile revient à l'état $0$~:
\[
A(t) = t\bigl(q\,A(t) + p\,B(t)\bigr),
\qquad
B(t) = t\bigl(p + q\,A(t)\bigr).
\]
En substituant~: $A(1 - qt) = pt\,B = pt(pt + qtA)$, donc
\[
G_T(t) = A(t) = \frac{p^2t^2}{1 - qt - pq\,t^2} .
\]
En $t = 1$ le dénominateur vaut $1 - q - pq = p(1 - q) = p^2$~:
$G_T(1) = 1$, le jeu se termine presque sûrement (comme
l'~\cref{exo:b2:proba:6} l'a montré par récurrence). Dérivation
logarithmique en $1$~: $\E(T) = 2 - \frac{D'(1)}{D(1)}$ avec
$D(t) = 1 - qt - pqt^2$, $D'(1) = -q - 2pq$~:
\[
\E(T) = 2 + \frac{q + 2pq}{p^2} = \frac{2p^2 + q + 2pq}{p^2}
= \frac{1 + p}{p^2},
\]
qui vaut $6$ pour $p = \frac12$.
\end{solution}

\begin{solution}{pb:b2:genfun:1}
\textbf{1.} Pour $t \in \intoo01$, la règle de la chaîne sur $G_n = G
\circ G_{n-1}$ donne $G_n'(t) =
G'\bigl(G_{n-1}(t)\bigr)G_{n-1}'(t)$. Quand $t \to 1^-$,
$G_{n-1}(t) \uparrow 1$, et $G'$ est croissante de
limite à gauche $m$ en $1$, donc le premier facteur tend vers $m$~; par
récurrence le second tend vers $m^{n-1}$. D'après le~\cref{thm:b2:genfun:moments}, $\E(Z_n) = G_n'(1^-) = m^n$.

\textbf{2.} En dérivant une fois de plus,
\[
G_n'' = G''(G_{n-1})\,(G_{n-1}')^2 +
G'(G_{n-1})\,G_{n-1}'',
\]
et en faisant $t \to 1^-$~: $a_n = G''(1)m^{2(n-1)} + m\,
a_{n-1}$ avec $a_n = G_n''(1)$, $a_1 = G''(1)$. Pour $m \neq
1$ on vérifie par récurrence que $a_n = G''(1)\,m^{n-1}
\frac{m^n - 1}{m - 1}$ (la récurrence ajoute $G''(1)m^{2n-2}$
à $m\cdot G''(1)m^{n-2}\frac{m^{n-1}-1}{m-1}$, et
$m^{n-1} + \frac{m^{n-1}-1}{m-1} = \frac{m^n - 1}{m-1}$)~;
pour $m = 1$, $a_n = a_{n-1} + G''(1) = n\,G''(1)$.

\textbf{3.} $V(Z_n) = a_n + m^n - m^{2n}$ et $G''(1) =
\sigma^2 + m^2 - m$. Pour $m \neq 1$, le morceau $(m^2 -
m)m^{n-1}\frac{m^n-1}{m-1} = m^n(m^n - 1)$ annule $m^n -
m^{2n}$ exactement, laissant $V(Z_n) =
\sigma^2m^{n-1}\frac{m^n-1}{m-1}$. Pour $m = 1$~: $V(Z_n) =
nG''(1) = n\sigma^2$.

\textbf{4.} $Z_n$ est une variable entière positive, donc
$\P(Z_n > 0) = \P(Z_n \geq 1) \leq \E(Z_n) = m^n$ par Markov
(\cref{thm:b2:randomvar:markov}). Pour $m < 1$ ceci décroît
géométriquement --- et de façon \omterm{def:b2:series:summable}{sommable}, donc Borel--Cantelli donne
même que seul un nombre fini de générations sont non vides, ce qui est
de nouveau l'extinction.

\textbf{5.} Cauchy--Schwarz~: $\E(Z_n)^2 = \E(Z_n\mathbf
1_{Z_n>0})^2 \leq \E(Z_n^2)\,\P(Z_n > 0)$. Avec la question 3
et $m < 1$~:
\[
\E(Z_n^2) = V(Z_n) + m^{2n}
\leq \frac{\sigma^2m^{n-1}}{1-m} + m^{2n},
\]
donc, en divisant $m^{2n}$ par cette borne et en simplifiant par
$m^n$,
\[
\P(Z_n > 0) \geq \frac{m^n}{\frac{\sigma^2}{m(1-m)} + m^n}
\geq \Bigl(\frac{\sigma^2}{m(1-m)} + 1\Bigr)^{-1}m^n ,
\]
en utilisant $m^n \leq 1$ au dénominateur. Avec la question 4~:
$\P(Z_n > 0) \asymp m^n$.

\textbf{6.} $G(t) = q\sum_k(pt)^k = \frac{q}{1 - pt}$, et $m
= G'(1) = \frac{pq}{(1-p)^2} = \frac pq$. Sous-critique pour $p
< \frac12$, critique pour $p = \frac12$, surcritique pour $p
> \frac12$.

\textbf{7.} $G(t) = t$ s'écrit $pt^2 - t + q = 0$, de racines
$\frac{1 \pm \abs{p - q}}{2p}$, c'est-à-dire $1$ et $\frac qp =
\frac1m$. La probabilité d'extinction est le plus petit point
fixe dans $\intcc01$ (\cref{thm:b2:genfun:extinction})~:
$q_{\mathrm{ext}} = 1$ si $m \leq 1$, et $\frac1m$ si $m >
1$.

\textbf{8.} Pour $m \neq 1$, avec $p = \frac m{m+1}$, $q =
\frac1{m+1}$~: si $q_n = \frac{m^n - 1}{m^{n+1} - 1}$, alors
\[
1 - p\,q_n = \frac{(m+1)(m^{n+1} - 1) - m(m^n - 1)}
{(m+1)(m^{n+1} - 1)} = \frac{m^{n+2} - 1}{(m+1)(m^{n+1} -
1)},
\]
donc $q_{n+1} = \frac{q}{1 - pq_n} = \frac{m^{n+1} -
1}{m^{n+2} - 1}$~; le cas de base $q_0 = 0$ est vérifié. Pour $m = 1$~:
$G(t) = \frac1{2 - t}$ et $q_{n+1} = \frac1{2 -
\frac{n}{n+1}} = \frac{n+1}{n+2}$, avec $q_0 = 0$.

\textbf{9.} $1 - q_n = \frac{m^n(m - 1)}{m^{n+1} - 1}$. Pour
$m < 1$ le dénominateur tend vers $-1$~: $1 - q_n \sim (1 -
m)\,m^n$. Pour $m > 1$~:
\[
q_{\mathrm{ext}} - q_n = \frac1m - \frac{m^n - 1}{m^{n+1} -
1} = \frac{m - 1}{m\,(m^{n+1} - 1)} \sim \frac{m -
1}{m^{2}}\;m^{-n} .
\]
Et $G'(t) = \frac{pq}{(1 - pt)^2}$ évaluée en $t = \frac
qp$ (où $1 - pt = 1 - q = p$) donne $G'(q_{\mathrm{ext}})
= \frac qp = \frac1m$~: le rapport observé $m^{-1}$ est exactement
la dérivée au point fixe attractif.

\textbf{10.} Pour $p = \frac12$~: $G''(t) = \frac{1/4}{(1 -
t/2)^3}$, donc $G''(1) = 2$ et $\sigma^2 = G''(1) + m - m^2 =
2$. La forme close donne $1 - q_n = \frac1{n+1}$~: la
probabilité de survie décroît comme $1/n$ --- trop lentement pour être
\omterm{def:b2:series:summable}{sommable}, contrairement à tout taux sous-critique.

\textbf{11.} Récurrence~: $G_1(t) = \frac1{2-t}$ correspond à la
formule pour $n = 1$, et
\[
G(G_n(t)) = \cfrac{1}{2 - \cfrac{n - (n-1)t}{n+1 - nt}}
= \frac{n + 1 - nt}{2(n+1) - 2nt - n + (n-1)t}
= \frac{n+1 - nt}{n + 2 - (n+1)t} .
\]
Puis
\[
\frac{G_n(t) - q_n}{1 - q_n}
= (n+1)\,\Bigl(\frac{n - (n-1)t}{n+1 - nt} -
\frac{n}{n+1}\Bigr)
= \frac{t}{n + 1 - nt}
= \frac{\frac{t}{n+1}}{1 - \frac{n}{n+1}t} ,
\]
la \omterm{ex:b2:powerseries:fibonacci}{fonction génératrice} de la \omterm{def:b2:randomvar:law}{loi} géométrique $\mathcal G\bigl(\frac1{n+1}
\bigr)$ sur $\N^*$ (\cref{ex:b2:genfun:classical})~: conditionnée à la
survie, $\P(Z_n = k \mid Z_n > 0) =
\frac1{n+1}\bigl(\frac n{n+1}\bigr)^{k-1}$, de moyenne
conditionnelle $n + 1$. La moyenne non conditionnelle $1 = \E(Z_n)$ est
le produit d'une probabilité de survie tendant vers zéro et d'une
taille conditionnelle croissant linéairement.

\textbf{12.} Si la lignée s'éteint à la génération $n$, alors
$Y = Z_0 + \dots + Z_{n-1}$ est finie~; si elle ne s'éteint
jamais, $Y \geq \sum_n 1 = \infty$. Donc $\{Y < \infty\}$ est
l'\omterm{def:b2:proba:space}{événement} d'extinction et $\P(Y < \infty) = q_{\mathrm{ext}}$.
La dérivation de $H(t) = tG(H(t))$ dans
l'~\cref{exo:b2:genfun:7} --- l'ancêtre contribue le
facteur $t$, ses enfants fondent des copies indépendantes de $Y$
comptées via $G$ --- n'a utilisé que le~\cref{thm:b2:genfun:compound}, valide dans tout régime.

\textbf{13.} Avec $G(s) = \frac{1 + s^2}2$ l'équation s'écrit
$tH^2 - 2H + t = 0$, donc $H = \frac{1 - \sqrt{1 - t^2}}{t}$
(la racine avec $H(0) = 0$). En comparant avec la série de Catalan
$C(x) = \frac{1 - \sqrt{1 - 4x}}{2x}$
(\cref{ex:b2:powerseries:catalan})~: $H(t) = \frac
t2\,C\bigl(\frac{t^2}4\bigr) =
\sum_{k\geq0}C_k\,\frac{t^{2k+1}}{2^{2k+1}}$, c'est-à-dire $\P(Y =
2k+1) = C_k2^{-2k-1}$. Vérifications~: $\P(Y = 1) = C_0/2 = \frac12$
(l'ancêtre n'a pas d'enfant)~; $\P(Y = 3) = C_1/8 = \frac18$
(deux enfants, tous deux sans enfant~: $\frac12\cdot\frac12\cdot
\frac12$).

\textbf{14.} En dérivant $H = tG(H)$ sur $\intoo01$ et
en faisant $t \to 1^-$ (limites monotones comme dans
le~\cref{thm:b2:genfun:moments})~: $H'(1)\bigl(1 - G'(H(1))\bigr)
= G(H(1))$. Dans le cas sous-critique $H(1) = 1$ et $\E(Y) =
H'(1) = \frac1{1 - m}$. Dans le cas critique $G'(1) = 1$
annule le facteur de gauche tandis que le membre de droite vaut $1$~:
aucun $H'(1)$ fini ne peut exister, donc $\E(Y) = \infty$ --- pourtant
$\P(Y < \infty) = q = 1$.

\textbf{15.} $C_k = \frac1{k+1}\binom{2k}k \sim
\frac{4^k}{\sqrt\pi\,k^{3/2}}$ d'après
l'~\cref{ex:b2:comparison:centralbinomial}, donc
\[
\P(Y = 2k+1) = \frac{C_k}{2\cdot4^{k}} \sim
\frac1{2\sqrt\pi\,k^{3/2}} .
\]
En sommant la queue (comparaison avec $\int_K^\infty
k^{-3/2}\dd k = 2K^{-1/2}$, par-dessus et par-dessous)~: $\P(Y > 2K)
\asymp K^{-1/2}$, c'est-à-dire $\P(Y > n) \asymp n^{-1/2}$ --- une
queue lourde d'\omterm{def:b2:randomvar:expectation}{espérance} infinie, quantifiant la question 14.

\textbf{16.} Les deux objets critiques --- le temps de retour
de la marche équilibrée (le problème du week-end du~\cref{ch:b2:proba}) et
la descendance totale critique --- sont presque sûrement finis
d'\omterm{def:b2:randomvar:expectation}{espérance} infinie, avec des \omterm{def:b2:randomvar:law}{lois} locales d'exposant $-3/2$ et des
queues d'exposant $-1/2$. Ce n'est pas une coïncidence~: explorer un
arbre généalogique enfant par enfant produit un chemin $\pm1$ (un pas
vers le haut par naissance, un pas vers le bas par mort) qui est
exactement une marche équilibrée, et $Y$ devient un temps de premier
passage. La criticité signifie dérive nulle~: le processus est
toujours au bord de l'extinction et de l'explosion, et les
fluctuations à l'échelle $\sqrt{}$ de l'aléa sans dérive produisent
précisément ces exposants.

\textbf{17.} $G''(t) = \sum_{n\geq2}n(n-1)p_nt^{n-2}$ a des
termes positifs, donc elle est croissante sur $\intco01$ de limite
finie $G''(1) = \sigma^2$ (la criticité fait
$\E Z_1(Z_1 - 1) = \sigma^2$)~; une fonction croissante de
limite égale à la valeur au bord est \omterm{def:b2:metric:continuity}{continue} en $1$.
Taylor avec reste intégral au point $1$~:
\[
G(t) = 1 + (t - 1) + \int_1^t(t - s)G''(s)\,\dd s
= t + \frac{G''(1)}2(1-t)^2 + o\bigl((1-t)^2\bigr),
\]
puisque $G''(s) = G''(1) + o(1)$ quand $s \to 1^-$.

\textbf{18.} En réduisant au même dénominateur, $h(t) =
\frac{G(t) - t}{(1 - G(t))(1 - t)}$. Par la question 17 le
numérateur est $b(1-t)^2 + o((1-t)^2)$ et $1 - G(t) = (1 -
t)\bigl(1 - b(1-t) + o(1-t)\bigr)$, donc $h(t) \to b$.

\textbf{19.} Par définition de $h$ en $t = q_j$ et $G(q_j) =
q_{j+1}$~: $\frac1{1 - q_{j+1}} - \frac1{1-q_j} = h(q_j)$~;
en sommant à partir de $j = 0$ ($q_0 = 0$) on obtient l'égalité
affichée. Puisque le processus critique s'éteint, $q_j \uparrow 1$,
donc $h(q_j) \to b$ et la moyenne de Cesàro $\frac1n\sum_{j<n}h(q_j)
\to b$~: $\frac1{1-q_n} \sim bn$, c'est-à-dire
\[
\P(Z_n > 0) \sim \frac1{bn} = \frac{2}{\sigma^2 n} .
\]

\textbf{20.} Cas géométrique critique~: $\sigma^2 = 2$
(question 10), donc Kolmogorov prédit $1 - q_n \sim \frac1n$
--- et la valeur \omterm{def:b2:multint:exact}{exacte} est $\frac1{n+1}$.

\textbf{21.} Puisque $Z_n\mathbf 1_{Z_n > 0} = Z_n$, $\E(Z_n
\mid Z_n > 0) = \frac{\E(Z_n)}{\P(Z_n > 0)} = \frac1{1 -
q_n} \sim \frac{\sigma^2n}2$. Dans le cas géométrique c'est
$n + 1$, correspondant exactement à la question 11 ($\sigma^2 = 2$). Le
tableau critique~: l'extinction est certaine, la taille moyenne est
figée à $1$, et les rares lignées survivantes ont une taille croissant
linéairement --- chaque facteur équilibrant l'autre.

\textbf{22.} Pour une reproduction $\mathcal P(\lambda)$, $G(t) =
\eu^{\lambda(t-1)}$ et la probabilité d'extinction est la
plus petite racine de $q = \eu^{\lambda(q-1)}$. Numériquement~:
$\lambda = 1.5$ donne $q \approx 0.417$ (itérer $q \mapsto
\eu^{1.5(q-1)}$~: $0, 0.223, 0.312, 0.356, \dots \to 0.4172$)~;
$\lambda = 2$ donne $q \approx 0.203$. Ainsi un unique cas
index déclenche une épidémie majeure avec probabilité $58\%$
($\lambda = 1.5$) ou $80\%$ ($\lambda = 2$) --- probable, non
certaine. L'itération partant de $q_0 = 0$ \omterm{def:b2:integration:improper}{converge} vers la
\emph{plus petite} racine car $G$ est croissante~: par
récurrence $q_n \leq r$ pour tout point fixe $r$, et $(q_n)$
croît (c'est $\P(Z_n = 0)$), donc sa limite est un point fixe
en dessous de tous les autres.

\textbf{23.} Les $k$ ancêtres fondent des arbres généalogiques
\omterm{def:b2:proba:independence}{indépendants}, et l'extinction totale est l'intersection de $k$
\omterm{def:b2:proba:space}{événements} d'extinction \omterm{def:b2:proba:independence}{indépendants}~: probabilité $q^k$. Pour
$\lambda = 1.5$~: la probabilité d'épidémie $1 - q^k \geq 0.99$
requiert $q^k \leq 0.01$, c'est-à-dire $k \geq
\frac{\ln 0.01}{\ln 0.417} \approx 5.3$~: six cas initiaux
rendent l'épidémie certaine à $99\%$.

\textbf{24.} \emph{$G'(q) < 1$~:} $G - \mathrm{id}$ est \omterm{def:b2:affine:convex}{convexe}
et s'annule en $q$ et $1$, donc elle est $\leq 0$ sur $\intcc
q1$~; si $G'(q) = 1$, la \omterm{def:b2:curves:arc}{tangente} en $q$ (que la convexité
place sous $G$) forcerait $G(t) \geq t$ sur $\intcc q1$,
d'où $G \equiv \mathrm{id}$ sur cet intervalle, annulant tous les
coefficients $p_n$ ($n \geq 2$) et contredisant $m > 1$.
\emph{Convergence géométrique~:} $q_n < q$ pour tout $n$ (récurrence,
$G$ croissante), et le théorème des accroissements finis donne
$q - q_{n+1} = G'(c_n)(q - q_n)$ avec $c_n \in \intoo{q_n}q$, donc
$G'(c_n) \leq G'(q) < 1$ et $q - q_n \leq q\,G'(q)^n$.
\emph{Processus compagnon~:} $\widehat G(t) = G(qt)/q =
\sum_kp_kq^{k-1}t^k$ a des coefficients positifs et
$\widehat G(1) = G(q)/q = 1$~: une \omterm{ex:b2:powerseries:fibonacci}{fonction génératrice}~; sa moyenne est
$\widehat G'(1) = G'(q) < 1$~: sous-critique. Famille géométrique~:
$G(t) = \frac{q}{1-pt}$, $q_{\mathrm{ext}} = \frac qp$, et
\[
\widehat G(t) = \frac pq\cdot\frac{q}{1 - p\frac qp t}
= \frac{p}{1 - qt} :
\]
la \omterm{def:b2:randomvar:law}{loi} de reproduction géométrique avec $p$ et $q$ échangés ---
le processus surcritique vu sur son \omterm{def:b2:proba:space}{événement} d'extinction est le
processus sous-critique miroir.

\textbf{25.} Le tableau~: $m < 1$~: $q = 1$, $\P(Z_n > 0)
\asymp m^n$ (questions 4--5), $\E(Y) = \frac1{1-m}$,
générations survivantes de moyenne conditionnelle bornée. $m = 1$~:
$q = 1$, $\P(Z_n > 0) \sim \frac2{\sigma^2n}$ (Kolmogorov),
$\E(Y) = \infty$ avec $\P(Y > n) \asymp n^{-1/2}$, survivants
de taille $\sim \frac{\sigma^2n}2$. $m > 1$~: $q < 1$ est le
plus petit point fixe, $q - q_n = O(G'(q)^n)$, croissance $\E(Z_n) = m^n$,
et conditionné à mourir le processus est le
compagnon sous-critique (question 24). Les outils~: la composition
des \omterm{ex:b2:powerseries:fibonacci}{fonctions génératrices} a transformé la récurrence de population en
itération de fonction~; la convexité a fixé la géométrie des points
fixes~; Taylor en $1^-$ a converti les hypothèses de moments en
développements locaux~; et la moyenne de Cesàro a extrait le $1/n$ de
Kolmogorov d'une somme télescopique. Le volume de troisième année
ajoute la martingale $Z_n/m^n$ --- dont la limite presque sûre raffine
$\E(Z_n) = m^n$ en un taux de croissance trajectoire par trajectoire
--- et le théorème de Yaglom, la \omterm{def:b2:randomvar:law}{loi} limite derrière la géométrie
conditionnelle observée à la question 11.
\end{solution}
